% GENERATED FILE. Edit canonical lessons, then run npm run generate:books.
% canonical-source-hash: fnv1a64:027bc966f4db9ef6
% canonical-lessons: ES-C596-temario, ES-C596-matricula, ES-C596-tutoria, ES-C596-facultad, ES-C596-master

\chapter{En la universidad}
\label{ch:es-en-la-universidad}
\hlchaptermodality{596}

\begin{chapteropening}
\textbf{By the end of this chapter:} \emph{I can talk about the syllabus, enrolment, tutorials and a master's degree.}

Complete the last lesson of chapter 596: I can talk about the syllabus, enrolment, tutorials and a master's degree.
\end{chapteropening}

\section[el temario]{el temario — the syllabus, the list of topics}

\label{lesson:ES-C596-temario}



\begin{quote}
Before the new one: say the Spanish for ambitious, then the Spanish for to succeed.
\end{quote}

\subsection*{You'll want to know: el temario}
\textbf{el temario} — \textquotedblleft{}the syllabus, the list of topics\textquotedblright{}.

\begin{quote}
El temario del examen es muy largo.
\end{quote}

\begin{quote}
\emph{The syllabus for the exam is very long.}
\end{quote}

In Spain it is the set of topics you study for an exam.

\begin{cousinweb}
\emph{temario} is built on \emph{tema} with \emph{-ario}; \emph{tema} comes from Latin \emph{thema}, taken from Greek \emph{thema} 'a thing laid down', from \emph{tithenai} 'to put' (EtymDB).
\end{cousinweb}

\begin{grammarlens}[title={What you've built}]
One more word for telling, explaining and describing at B1.
\end{grammarlens}

\subsection*{Your turn}
\begin{itemize}
\raggedright
  \item \emph{Say it:} \emph{el temario}
  \item \emph{Say it:} \emph{el temario}, once more
  \item \emph{Say it:} a sentence of your own with \emph{el temario}
  \item \emph{Recall:} say the Spanish for ambitious, then the Spanish for to succeed, then say \emph{el temario} again
\end{itemize}

\begin{tcolorbox}[breakable,colback=teal!4,colframe=teal!35!black,title={Before you move on}]
What does el temario mean? (\textquotedblleft{}The syllabus, the list of topics\textquotedblright{}.) Say it once more.
\end{tcolorbox}
\section[la matrícula]{la matrícula — the enrolment, the registration; the number plate}

\label{lesson:ES-C596-matricula}



\begin{quote}
Before the new one: say the Spanish for to succeed, then the Spanish for the syllabus.
\end{quote}

\subsection*{You'll want to know: la matrícula}
\textbf{la matrícula} — \textquotedblleft{}the enrolment, the registration; the number plate\textquotedblright{}.

\begin{quote}
Pagué trescientos euros de matrícula.
\end{quote}

\begin{quote}
\emph{I paid three hundred euros for enrolment.}
\end{quote}

In Spain it is also the car number plate; in much of Latin America that is \emph{la placa}.

\begin{cousinweb}
\emph{matrícula} is a learned borrowing from Late Latin \emph{matricula} 'a list, a register', a small form of \emph{matrix}, built on \emph{mater} 'mother' (EtymDB).
\end{cousinweb}

\begin{grammarlens}[title={What you've built}]
One more word for telling, explaining and describing at B1.
\end{grammarlens}

\subsection*{Your turn}
\begin{itemize}
\raggedright
  \item \emph{Say it:} \emph{la matrícula}
  \item \emph{Say it:} \emph{la matrícula}, once more
  \item \emph{Say it:} a sentence of your own with \emph{la matrícula}
  \item \emph{Recall:} say the Spanish for to succeed, then the Spanish for the syllabus, then say \emph{la matrícula} again
\end{itemize}

\begin{tcolorbox}[breakable,colback=teal!4,colframe=teal!35!black,title={Before you move on}]
What does la matrícula mean? (\textquotedblleft{}The enrolment, the registration; the number plate\textquotedblright{}.) Say it once more.
\end{tcolorbox}
\section[la tutoría]{la tutoría — the tutorial, the office hours}

\label{lesson:ES-C596-tutoria}



\begin{quote}
Before the new one: say the Spanish for the syllabus, then the Spanish for the enrolment.
\end{quote}

\subsection*{You'll want to know: la tutoría}
\textbf{la tutoría} — \textquotedblleft{}the tutorial, the office hours\textquotedblright{}.

\begin{quote}
Tengo tutoría con la profesora el jueves.
\end{quote}

\begin{quote}
\emph{I have a tutorial with the teacher on Thursday.}
\end{quote}

The time when a teacher sees students one by one.

\begin{cousinweb}
\emph{tutoría} is built on \emph{tutor}, a learned borrowing from Latin \emph{tutor} 'guardian, protector' (EtymDB), from \emph{tueri} 'to watch over'.
\end{cousinweb}

\begin{grammarlens}[title={What you've built}]
One more word for telling, explaining and describing at B1.
\end{grammarlens}

\subsection*{Your turn}
\begin{itemize}
\raggedright
  \item \emph{Say it:} \emph{la tutoría}
  \item \emph{Say it:} \emph{la tutoría}, once more
  \item \emph{Say it:} a sentence of your own with \emph{la tutoría}
  \item \emph{Recall:} say the Spanish for the syllabus, then the Spanish for the enrolment, then say \emph{la tutoría} again
\end{itemize}

\begin{tcolorbox}[breakable,colback=teal!4,colframe=teal!35!black,title={Before you move on}]
What does la tutoría mean? (\textquotedblleft{}The tutorial, the office hours\textquotedblright{}.) Say it once more.
\end{tcolorbox}
\section[la facultad]{la facultad — the faculty, the school of a university; ability}

\label{lesson:ES-C596-facultad}



\begin{quote}
Before the new one: say the Spanish for the enrolment, then the Spanish for the tutorial.
\end{quote}

\subsection*{You'll want to know: la facultad}
\textbf{la facultad} — \textquotedblleft{}the faculty, the school of a university; ability\textquotedblright{}.

\begin{quote}
Estudio en la Facultad de Derecho.
\end{quote}

\begin{quote}
\emph{I study at the Faculty of Law.}
\end{quote}

\emph{Tener facultades para} — to have the ability for.

\begin{cousinweb}
\emph{facultad} is a learned borrowing from Latin \emph{facultas} 'skill, ability', built on an old form of \emph{facilis} 'easy to do' (EtymDB).
\end{cousinweb}

\begin{grammarlens}[title={What you've built}]
One more word for telling, explaining and describing at B1.
\end{grammarlens}

\subsection*{Your turn}
\begin{itemize}
\raggedright
  \item \emph{Say it:} \emph{la facultad}
  \item \emph{Say it:} \emph{la facultad}, once more
  \item \emph{Say it:} a sentence of your own with \emph{la facultad}
  \item \emph{Recall:} say the Spanish for the enrolment, then the Spanish for the tutorial, then say \emph{la facultad} again
\end{itemize}

\begin{tcolorbox}[breakable,colback=teal!4,colframe=teal!35!black,title={Before you move on}]
What does la facultad mean? (\textquotedblleft{}The faculty, the school of a university; ability\textquotedblright{}.) Say it once more.
\end{tcolorbox}
\section[el máster]{el máster — the master's degree}

\label{lesson:ES-C596-master}



\begin{quote}
Before the new one: say the Spanish for the tutorial, then the Spanish for the faculty.
\end{quote}

\subsection*{You'll want to know: el máster}
\textbf{el máster} — \textquotedblleft{}the master's degree\textquotedblright{}.

\begin{quote}
Quiero hacer un máster en historia.
\end{quote}

\begin{quote}
\emph{I want to do a master's in history.}
\end{quote}

This is the word in Spain; much of Latin America says \emph{la maestría}.

\begin{cousinweb}
\emph{máster} was borrowed from English \emph{master}, which comes from Latin \emph{magister} 'master, teacher' (EtymDB).
\end{cousinweb}

\begin{grammarlens}[title={What you've built}]
One more word for telling, explaining and describing at B1.
\end{grammarlens}

\subsection*{Your turn}
\begin{itemize}
\raggedright
  \item \emph{Say it:} \emph{el máster}
  \item \emph{Say it:} \emph{el máster}, once more
  \item \emph{Say it:} a sentence of your own with \emph{el máster}
  \item \emph{Recall:} say the Spanish for the tutorial, then the Spanish for the faculty, then say \emph{el máster} again
\end{itemize}

\begin{tcolorbox}[breakable,colback=teal!4,colframe=teal!35!black,title={Before you move on}]
What does el máster mean? (\textquotedblleft{}The master's degree\textquotedblright{}.) Say it once more.
\end{tcolorbox}
